\documentclass[11pt,letterpaper]{article}

\usepackage{fontspec}
\usepackage[spanish,es-noshorthands]{babel}
\usepackage[margin=2.5cm]{geometry}
\usepackage{amsmath}
\usepackage{xcolor}
\usepackage{luacolor}
\usepackage{lua-ul}
\usepackage[hidelinks]{hyperref}

% -------------------------------------------------------------
\definecolor{cFoster}{HTML}{FFE97A}
\definecolor{cSalak}{HTML}{B8EBB0}
\definecolor{cOussidi}{HTML}{AEDBFF}
\definecolor{cHe}{HTML}{FFC2E2}

\newcommand{\fos}[1]{\highLight[cFoster]{#1}}
\newcommand{\sal}[1]{\highLight[cSalak]{#1}}
\newcommand{\ous}[1]{\highLight[cOussidi]{#1}}
\newcommand{\he}[1]{\highLight[cHe]{#1}}

\setlength{\parindent}{0pt}
\setlength{\parskip}{0.5em}

\title{¿Qué es la IA generativa?}
\author{}
\date{}

\begin{document}
\maketitle

\fos{Foster (2019)} \quad \sal{Salakhutdinov (2015)} \quad \ous{Oussidi y Elhassouny (2018)} \quad \he{He, Cao y Tan (2025)}

% ----------
\section{Definición}

\he{La IA generativa es un conjunto de algoritmos y modelos capaces de producir contenido nuevo, como texto, imágenes, video, código o estrategias, con una creatividad y adaptabilidad parecidas a las humanas (He et al., 2025).}

\fos{Un modelo generativo describe cómo se genera un conjunto de datos mediante un modelo probabilístico, al muestrear de él se obtienen datos nuevos (Foster, 2019).} \ous{Es decir, busca capturar la distribución que genera una clase de datos para producir datos similares (Oussidi y Elhassouny, 2018).}

% ----------
\section{Idea principal}

\fos{Se parte de observaciones $\mathbf{X}$ generadas por una distribución desconocida $p_{\text{data}}$. El modelo $p_{\text{model}}$ intenta imitarla, y es bueno si sus muestras parecen reales y no son copias de los datos de entrenamiento (Foster, 2019).}

\fos{A diferencia de un modelo discriminativo, que estima $p(y \mid \mathbf{x})$, un modelo generativo estima $p(\mathbf{x})$, por eso suele entrenarse con datos sin etiquetar (Foster, 2019).}

% ----------
\section{Por qué el aprendizaje profundo}

\sal{Las arquitecturas profundas, con muchas capas no lineales y variables latentes, pueden aprender representaciones complejas a partir de grandes cantidades de datos sin etiquetar (Salakhutdinov, 2015).} \fos{La clave es el espacio latente: describir cada dato con pocas dimensiones y aprender a mapearlas de vuelta al espacio original (Foster, 2019).}

% ----------
\section{Los modelos principales}

\begin{itemize}
  \item \sal{Máquinas de Boltzmann y redes de creencia profunda: modelos basados en energía con capas ocultas (Salakhutdinov, 2015).}
  \item \ous{Autocodificadores variacionales (VAE): codifican los datos en un espacio latente gaussiano del que se puede muestrear (Oussidi y Elhassouny, 2018).}
  \item \fos{Redes generativas adversarias (GAN): un generador y un discriminador compiten entre sí (Foster, 2019).}
  \item \ous{Modelos autorregresivos: generan un elemento a la vez a partir de los anteriores (Oussidi y Elhassouny, 2018).}
  \item \he{Modelos de difusión: aprenden a convertir ruido en datos (He et al., 2025).}
  \item \he{Transformers y modelos fundacionales: entrenados a gran escala y adaptables a muchas tareas, como GPT (He et al., 2025).}
\end{itemize}

% ----------
\section{Conclusión}

La IA generativa \he{crea contenido nuevo} al \fos{aprender cómo son los datos y generar ejemplos parecidos}. No copia, aprende sus patrones. \he{Sus retos incluyen las alucinaciones, el costo y el uso indebido (He et al., 2025).}

% ----------
\section*{Referencias}

\fos{Foster, D. (2019). \emph{Generative Deep Learning}. O'Reilly Media.}

\sal{Salakhutdinov, R. (2015). Learning Deep Generative Models. \emph{Annual Review of Statistics and Its Application}, 2, 361--385.}

\ous{Oussidi, A. y Elhassouny, A. (2018). Deep Generative Models: Survey. \emph{ISCV 2018}, IEEE.}

\he{He, R., Cao, J. y Tan, T. (2025). Generative Artificial Intelligence: A Historical Perspective. \emph{National Science Review}, 12(5).}

\end{document}
